\documentclass[11pt]{article}

    \usepackage[breakable]{tcolorbox}
    \usepackage{parskip} % Stop auto-indenting (to mimic markdown behaviour)
    

    % Basic figure setup, for now with no caption control since it's done
    % automatically by Pandoc (which extracts ![](path) syntax from Markdown).
    \usepackage{graphicx}
    % Keep aspect ratio if custom image width or height is specified
    \setkeys{Gin}{keepaspectratio}
    % Maintain compatibility with old templates. Remove in nbconvert 6.0
    \let\Oldincludegraphics\includegraphics
    % Ensure that by default, figures have no caption (until we provide a
    % proper Figure object with a Caption API and a way to capture that
    % in the conversion process - todo).
    \usepackage{caption}
    \DeclareCaptionFormat{nocaption}{}
    \captionsetup{format=nocaption,aboveskip=0pt,belowskip=0pt}

    \usepackage{float}
    \floatplacement{figure}{H} % forces figures to be placed at the correct location
    \usepackage{xcolor} % Allow colors to be defined
    \usepackage{enumerate} % Needed for markdown enumerations to work
    \usepackage{geometry} % Used to adjust the document margins
    \usepackage{amsmath} % Equations
    \usepackage{amssymb} % Equations
    \usepackage{textcomp} % defines textquotesingle
    % Hack from http://tex.stackexchange.com/a/47451/13684:
    \AtBeginDocument{%
        \def\PYZsq{\textquotesingle}% Upright quotes in Pygmentized code
    }
    \usepackage{upquote} % Upright quotes for verbatim code
    \usepackage{eurosym} % defines \euro

    \usepackage{iftex}
    \ifPDFTeX
        \usepackage[T1]{fontenc}
        \IfFileExists{alphabeta.sty}{
              \usepackage{alphabeta}
          }{
              \usepackage[mathletters]{ucs}
              \usepackage[utf8x]{inputenc}
          }
    \else
        \usepackage{fontspec}
        \usepackage{unicode-math}
    \fi

    \usepackage{fancyvrb} % verbatim replacement that allows latex
    \usepackage{grffile} % extends the file name processing of package graphics
                         % to support a larger range
    \makeatletter % fix for old versions of grffile with XeLaTeX
    \@ifpackagelater{grffile}{2019/11/01}
    {
      % Do nothing on new versions
    }
    {
      \def\Gread@@xetex#1{%
        \IfFileExists{"\Gin@base".bb}%
        {\Gread@eps{\Gin@base.bb}}%
        {\Gread@@xetex@aux#1}%
      }
    }
    \makeatother
    \usepackage[Export]{adjustbox} % Used to constrain images to a maximum size
    \adjustboxset{max size={0.9\linewidth}{0.9\paperheight}}

    % The hyperref package gives us a pdf with properly built
    % internal navigation ('pdf bookmarks' for the table of contents,
    % internal cross-reference links, web links for URLs, etc.)
    \usepackage{hyperref}
    % The default LaTeX title has an obnoxious amount of whitespace. By default,
    % titling removes some of it. It also provides customization options.
    \usepackage{titling}
    \usepackage{longtable} % longtable support required by pandoc >1.10
    \usepackage{booktabs}  % table support for pandoc > 1.12.2
    \usepackage{array}     % table support for pandoc >= 2.11.3
    \usepackage{calc}      % table minipage width calculation for pandoc >= 2.11.1
    \usepackage[inline]{enumitem} % IRkernel/repr support (it uses the enumerate* environment)
    \usepackage[normalem]{ulem} % ulem is needed to support strikethroughs (\sout)
                                % normalem makes italics be italics, not underlines
    \usepackage{soul}      % strikethrough (\st) support for pandoc >= 3.0.0
    \usepackage{mathrsfs}
    

    
    % Colors for the hyperref package
    \definecolor{urlcolor}{rgb}{0,.145,.698}
    \definecolor{linkcolor}{rgb}{.71,0.21,0.01}
    \definecolor{citecolor}{rgb}{.12,.54,.11}

    % ANSI colors
    \definecolor{ansi-black}{HTML}{3E424D}
    \definecolor{ansi-black-intense}{HTML}{282C36}
    \definecolor{ansi-red}{HTML}{E75C58}
    \definecolor{ansi-red-intense}{HTML}{B22B31}
    \definecolor{ansi-green}{HTML}{00A250}
    \definecolor{ansi-green-intense}{HTML}{007427}
    \definecolor{ansi-yellow}{HTML}{DDB62B}
    \definecolor{ansi-yellow-intense}{HTML}{B27D12}
    \definecolor{ansi-blue}{HTML}{208FFB}
    \definecolor{ansi-blue-intense}{HTML}{0065CA}
    \definecolor{ansi-magenta}{HTML}{D160C4}
    \definecolor{ansi-magenta-intense}{HTML}{A03196}
    \definecolor{ansi-cyan}{HTML}{60C6C8}
    \definecolor{ansi-cyan-intense}{HTML}{258F8F}
    \definecolor{ansi-white}{HTML}{C5C1B4}
    \definecolor{ansi-white-intense}{HTML}{A1A6B2}
    \definecolor{ansi-default-inverse-fg}{HTML}{FFFFFF}
    \definecolor{ansi-default-inverse-bg}{HTML}{000000}

    % common color for the border for error outputs.
    \definecolor{outerrorbackground}{HTML}{FFDFDF}

    % commands and environments needed by pandoc snippets
    % extracted from the output of `pandoc -s`
    \providecommand{\tightlist}{%
      \setlength{\itemsep}{0pt}\setlength{\parskip}{0pt}}
    \DefineVerbatimEnvironment{Highlighting}{Verbatim}{commandchars=\\\{\}}
    % Add ',fontsize=\small' for more characters per line
    \newenvironment{Shaded}{}{}
    \newcommand{\KeywordTok}[1]{\textcolor[rgb]{0.00,0.44,0.13}{\textbf{{#1}}}}
    \newcommand{\DataTypeTok}[1]{\textcolor[rgb]{0.56,0.13,0.00}{{#1}}}
    \newcommand{\DecValTok}[1]{\textcolor[rgb]{0.25,0.63,0.44}{{#1}}}
    \newcommand{\BaseNTok}[1]{\textcolor[rgb]{0.25,0.63,0.44}{{#1}}}
    \newcommand{\FloatTok}[1]{\textcolor[rgb]{0.25,0.63,0.44}{{#1}}}
    \newcommand{\CharTok}[1]{\textcolor[rgb]{0.25,0.44,0.63}{{#1}}}
    \newcommand{\StringTok}[1]{\textcolor[rgb]{0.25,0.44,0.63}{{#1}}}
    \newcommand{\CommentTok}[1]{\textcolor[rgb]{0.38,0.63,0.69}{\textit{{#1}}}}
    \newcommand{\OtherTok}[1]{\textcolor[rgb]{0.00,0.44,0.13}{{#1}}}
    \newcommand{\AlertTok}[1]{\textcolor[rgb]{1.00,0.00,0.00}{\textbf{{#1}}}}
    \newcommand{\FunctionTok}[1]{\textcolor[rgb]{0.02,0.16,0.49}{{#1}}}
    \newcommand{\RegionMarkerTok}[1]{{#1}}
    \newcommand{\ErrorTok}[1]{\textcolor[rgb]{1.00,0.00,0.00}{\textbf{{#1}}}}
    \newcommand{\NormalTok}[1]{{#1}}

    % Additional commands for more recent versions of Pandoc
    \newcommand{\ConstantTok}[1]{\textcolor[rgb]{0.53,0.00,0.00}{{#1}}}
    \newcommand{\SpecialCharTok}[1]{\textcolor[rgb]{0.25,0.44,0.63}{{#1}}}
    \newcommand{\VerbatimStringTok}[1]{\textcolor[rgb]{0.25,0.44,0.63}{{#1}}}
    \newcommand{\SpecialStringTok}[1]{\textcolor[rgb]{0.73,0.40,0.53}{{#1}}}
    \newcommand{\ImportTok}[1]{{#1}}
    \newcommand{\DocumentationTok}[1]{\textcolor[rgb]{0.73,0.13,0.13}{\textit{{#1}}}}
    \newcommand{\AnnotationTok}[1]{\textcolor[rgb]{0.38,0.63,0.69}{\textbf{\textit{{#1}}}}}
    \newcommand{\CommentVarTok}[1]{\textcolor[rgb]{0.38,0.63,0.69}{\textbf{\textit{{#1}}}}}
    \newcommand{\VariableTok}[1]{\textcolor[rgb]{0.10,0.09,0.49}{{#1}}}
    \newcommand{\ControlFlowTok}[1]{\textcolor[rgb]{0.00,0.44,0.13}{\textbf{{#1}}}}
    \newcommand{\OperatorTok}[1]{\textcolor[rgb]{0.40,0.40,0.40}{{#1}}}
    \newcommand{\BuiltInTok}[1]{{#1}}
    \newcommand{\ExtensionTok}[1]{{#1}}
    \newcommand{\PreprocessorTok}[1]{\textcolor[rgb]{0.74,0.48,0.00}{{#1}}}
    \newcommand{\AttributeTok}[1]{\textcolor[rgb]{0.49,0.56,0.16}{{#1}}}
    \newcommand{\InformationTok}[1]{\textcolor[rgb]{0.38,0.63,0.69}{\textbf{\textit{{#1}}}}}
    \newcommand{\WarningTok}[1]{\textcolor[rgb]{0.38,0.63,0.69}{\textbf{\textit{{#1}}}}}
    \makeatletter
    \newsavebox\pandoc@box
    \newcommand*\pandocbounded[1]{%
      \sbox\pandoc@box{#1}%
      % scaling factors for width and height
      \Gscale@div\@tempa\textheight{\dimexpr\ht\pandoc@box+\dp\pandoc@box\relax}%
      \Gscale@div\@tempb\linewidth{\wd\pandoc@box}%
      % select the smaller of both
      \ifdim\@tempb\p@<\@tempa\p@
        \let\@tempa\@tempb
      \fi
      % scaling accordingly (\@tempa < 1)
      \ifdim\@tempa\p@<\p@
        \scalebox{\@tempa}{\usebox\pandoc@box}%
      % scaling not needed, use as it is
      \else
        \usebox{\pandoc@box}%
      \fi
    }
    \makeatother

    % Define a nice break command that doesn't care if a line doesn't already
    % exist.
    \def\br{\hspace*{\fill} \\* }
    % Math Jax compatibility definitions
    \def\gt{>}
    \def\lt{<}
    \let\Oldtex\TeX
    \let\Oldlatex\LaTeX
    \renewcommand{\TeX}{\textrm{\Oldtex}}
    \renewcommand{\LaTeX}{\textrm{\Oldlatex}}
    % Document parameters
    % Document title
    \title{assignments}
    
    
    
    
    
    
    
% Pygments definitions
\makeatletter
\def\PY@reset{\let\PY@it=\relax \let\PY@bf=\relax%
    \let\PY@ul=\relax \let\PY@tc=\relax%
    \let\PY@bc=\relax \let\PY@ff=\relax}
\def\PY@tok#1{\csname PY@tok@#1\endcsname}
\def\PY@toks#1+{\ifx\relax#1\empty\else%
    \PY@tok{#1}\expandafter\PY@toks\fi}
\def\PY@do#1{\PY@bc{\PY@tc{\PY@ul{%
    \PY@it{\PY@bf{\PY@ff{#1}}}}}}}
\def\PY#1#2{\PY@reset\PY@toks#1+\relax+\PY@do{#2}}

\@namedef{PY@tok@w}{\def\PY@tc##1{\textcolor[rgb]{0.73,0.73,0.73}{##1}}}
\@namedef{PY@tok@c}{\let\PY@it=\textit\def\PY@tc##1{\textcolor[rgb]{0.24,0.48,0.48}{##1}}}
\@namedef{PY@tok@cp}{\def\PY@tc##1{\textcolor[rgb]{0.61,0.40,0.00}{##1}}}
\@namedef{PY@tok@k}{\let\PY@bf=\textbf\def\PY@tc##1{\textcolor[rgb]{0.00,0.50,0.00}{##1}}}
\@namedef{PY@tok@kp}{\def\PY@tc##1{\textcolor[rgb]{0.00,0.50,0.00}{##1}}}
\@namedef{PY@tok@kt}{\def\PY@tc##1{\textcolor[rgb]{0.69,0.00,0.25}{##1}}}
\@namedef{PY@tok@o}{\def\PY@tc##1{\textcolor[rgb]{0.40,0.40,0.40}{##1}}}
\@namedef{PY@tok@ow}{\let\PY@bf=\textbf\def\PY@tc##1{\textcolor[rgb]{0.67,0.13,1.00}{##1}}}
\@namedef{PY@tok@nb}{\def\PY@tc##1{\textcolor[rgb]{0.00,0.50,0.00}{##1}}}
\@namedef{PY@tok@nf}{\def\PY@tc##1{\textcolor[rgb]{0.00,0.00,1.00}{##1}}}
\@namedef{PY@tok@nc}{\let\PY@bf=\textbf\def\PY@tc##1{\textcolor[rgb]{0.00,0.00,1.00}{##1}}}
\@namedef{PY@tok@nn}{\let\PY@bf=\textbf\def\PY@tc##1{\textcolor[rgb]{0.00,0.00,1.00}{##1}}}
\@namedef{PY@tok@ne}{\let\PY@bf=\textbf\def\PY@tc##1{\textcolor[rgb]{0.80,0.25,0.22}{##1}}}
\@namedef{PY@tok@nv}{\def\PY@tc##1{\textcolor[rgb]{0.10,0.09,0.49}{##1}}}
\@namedef{PY@tok@no}{\def\PY@tc##1{\textcolor[rgb]{0.53,0.00,0.00}{##1}}}
\@namedef{PY@tok@nl}{\def\PY@tc##1{\textcolor[rgb]{0.46,0.46,0.00}{##1}}}
\@namedef{PY@tok@ni}{\let\PY@bf=\textbf\def\PY@tc##1{\textcolor[rgb]{0.44,0.44,0.44}{##1}}}
\@namedef{PY@tok@na}{\def\PY@tc##1{\textcolor[rgb]{0.41,0.47,0.13}{##1}}}
\@namedef{PY@tok@nt}{\let\PY@bf=\textbf\def\PY@tc##1{\textcolor[rgb]{0.00,0.50,0.00}{##1}}}
\@namedef{PY@tok@nd}{\def\PY@tc##1{\textcolor[rgb]{0.67,0.13,1.00}{##1}}}
\@namedef{PY@tok@s}{\def\PY@tc##1{\textcolor[rgb]{0.73,0.13,0.13}{##1}}}
\@namedef{PY@tok@sd}{\let\PY@it=\textit\def\PY@tc##1{\textcolor[rgb]{0.73,0.13,0.13}{##1}}}
\@namedef{PY@tok@si}{\let\PY@bf=\textbf\def\PY@tc##1{\textcolor[rgb]{0.64,0.35,0.47}{##1}}}
\@namedef{PY@tok@se}{\let\PY@bf=\textbf\def\PY@tc##1{\textcolor[rgb]{0.67,0.36,0.12}{##1}}}
\@namedef{PY@tok@sr}{\def\PY@tc##1{\textcolor[rgb]{0.64,0.35,0.47}{##1}}}
\@namedef{PY@tok@ss}{\def\PY@tc##1{\textcolor[rgb]{0.10,0.09,0.49}{##1}}}
\@namedef{PY@tok@sx}{\def\PY@tc##1{\textcolor[rgb]{0.00,0.50,0.00}{##1}}}
\@namedef{PY@tok@m}{\def\PY@tc##1{\textcolor[rgb]{0.40,0.40,0.40}{##1}}}
\@namedef{PY@tok@gh}{\let\PY@bf=\textbf\def\PY@tc##1{\textcolor[rgb]{0.00,0.00,0.50}{##1}}}
\@namedef{PY@tok@gu}{\let\PY@bf=\textbf\def\PY@tc##1{\textcolor[rgb]{0.50,0.00,0.50}{##1}}}
\@namedef{PY@tok@gd}{\def\PY@tc##1{\textcolor[rgb]{0.63,0.00,0.00}{##1}}}
\@namedef{PY@tok@gi}{\def\PY@tc##1{\textcolor[rgb]{0.00,0.52,0.00}{##1}}}
\@namedef{PY@tok@gr}{\def\PY@tc##1{\textcolor[rgb]{0.89,0.00,0.00}{##1}}}
\@namedef{PY@tok@ge}{\let\PY@it=\textit}
\@namedef{PY@tok@gs}{\let\PY@bf=\textbf}
\@namedef{PY@tok@ges}{\let\PY@bf=\textbf\let\PY@it=\textit}
\@namedef{PY@tok@gp}{\let\PY@bf=\textbf\def\PY@tc##1{\textcolor[rgb]{0.00,0.00,0.50}{##1}}}
\@namedef{PY@tok@go}{\def\PY@tc##1{\textcolor[rgb]{0.44,0.44,0.44}{##1}}}
\@namedef{PY@tok@gt}{\def\PY@tc##1{\textcolor[rgb]{0.00,0.27,0.87}{##1}}}
\@namedef{PY@tok@err}{\def\PY@bc##1{{\setlength{\fboxsep}{\string -\fboxrule}\fcolorbox[rgb]{1.00,0.00,0.00}{1,1,1}{\strut ##1}}}}
\@namedef{PY@tok@kc}{\let\PY@bf=\textbf\def\PY@tc##1{\textcolor[rgb]{0.00,0.50,0.00}{##1}}}
\@namedef{PY@tok@kd}{\let\PY@bf=\textbf\def\PY@tc##1{\textcolor[rgb]{0.00,0.50,0.00}{##1}}}
\@namedef{PY@tok@kn}{\let\PY@bf=\textbf\def\PY@tc##1{\textcolor[rgb]{0.00,0.50,0.00}{##1}}}
\@namedef{PY@tok@kr}{\let\PY@bf=\textbf\def\PY@tc##1{\textcolor[rgb]{0.00,0.50,0.00}{##1}}}
\@namedef{PY@tok@bp}{\def\PY@tc##1{\textcolor[rgb]{0.00,0.50,0.00}{##1}}}
\@namedef{PY@tok@fm}{\def\PY@tc##1{\textcolor[rgb]{0.00,0.00,1.00}{##1}}}
\@namedef{PY@tok@vc}{\def\PY@tc##1{\textcolor[rgb]{0.10,0.09,0.49}{##1}}}
\@namedef{PY@tok@vg}{\def\PY@tc##1{\textcolor[rgb]{0.10,0.09,0.49}{##1}}}
\@namedef{PY@tok@vi}{\def\PY@tc##1{\textcolor[rgb]{0.10,0.09,0.49}{##1}}}
\@namedef{PY@tok@vm}{\def\PY@tc##1{\textcolor[rgb]{0.10,0.09,0.49}{##1}}}
\@namedef{PY@tok@sa}{\def\PY@tc##1{\textcolor[rgb]{0.73,0.13,0.13}{##1}}}
\@namedef{PY@tok@sb}{\def\PY@tc##1{\textcolor[rgb]{0.73,0.13,0.13}{##1}}}
\@namedef{PY@tok@sc}{\def\PY@tc##1{\textcolor[rgb]{0.73,0.13,0.13}{##1}}}
\@namedef{PY@tok@dl}{\def\PY@tc##1{\textcolor[rgb]{0.73,0.13,0.13}{##1}}}
\@namedef{PY@tok@s2}{\def\PY@tc##1{\textcolor[rgb]{0.73,0.13,0.13}{##1}}}
\@namedef{PY@tok@sh}{\def\PY@tc##1{\textcolor[rgb]{0.73,0.13,0.13}{##1}}}
\@namedef{PY@tok@s1}{\def\PY@tc##1{\textcolor[rgb]{0.73,0.13,0.13}{##1}}}
\@namedef{PY@tok@mb}{\def\PY@tc##1{\textcolor[rgb]{0.40,0.40,0.40}{##1}}}
\@namedef{PY@tok@mf}{\def\PY@tc##1{\textcolor[rgb]{0.40,0.40,0.40}{##1}}}
\@namedef{PY@tok@mh}{\def\PY@tc##1{\textcolor[rgb]{0.40,0.40,0.40}{##1}}}
\@namedef{PY@tok@mi}{\def\PY@tc##1{\textcolor[rgb]{0.40,0.40,0.40}{##1}}}
\@namedef{PY@tok@il}{\def\PY@tc##1{\textcolor[rgb]{0.40,0.40,0.40}{##1}}}
\@namedef{PY@tok@mo}{\def\PY@tc##1{\textcolor[rgb]{0.40,0.40,0.40}{##1}}}
\@namedef{PY@tok@ch}{\let\PY@it=\textit\def\PY@tc##1{\textcolor[rgb]{0.24,0.48,0.48}{##1}}}
\@namedef{PY@tok@cm}{\let\PY@it=\textit\def\PY@tc##1{\textcolor[rgb]{0.24,0.48,0.48}{##1}}}
\@namedef{PY@tok@cpf}{\let\PY@it=\textit\def\PY@tc##1{\textcolor[rgb]{0.24,0.48,0.48}{##1}}}
\@namedef{PY@tok@c1}{\let\PY@it=\textit\def\PY@tc##1{\textcolor[rgb]{0.24,0.48,0.48}{##1}}}
\@namedef{PY@tok@cs}{\let\PY@it=\textit\def\PY@tc##1{\textcolor[rgb]{0.24,0.48,0.48}{##1}}}

\def\PYZbs{\char`\\}
\def\PYZus{\char`\_}
\def\PYZob{\char`\{}
\def\PYZcb{\char`\}}
\def\PYZca{\char`\^}
\def\PYZam{\char`\&}
\def\PYZlt{\char`\<}
\def\PYZgt{\char`\>}
\def\PYZsh{\char`\#}
\def\PYZpc{\char`\%}
\def\PYZdl{\char`\$}
\def\PYZhy{\char`\-}
\def\PYZsq{\char`\'}
\def\PYZdq{\char`\"}
\def\PYZti{\char`\~}
% for compatibility with earlier versions
\def\PYZat{@}
\def\PYZlb{[}
\def\PYZrb{]}
\makeatother


    % For linebreaks inside Verbatim environment from package fancyvrb.
    \makeatletter
        \newbox\Wrappedcontinuationbox
        \newbox\Wrappedvisiblespacebox
        \newcommand*\Wrappedvisiblespace {\textcolor{red}{\textvisiblespace}}
        \newcommand*\Wrappedcontinuationsymbol {\textcolor{red}{\llap{\tiny$\m@th\hookrightarrow$}}}
        \newcommand*\Wrappedcontinuationindent {3ex }
        \newcommand*\Wrappedafterbreak {\kern\Wrappedcontinuationindent\copy\Wrappedcontinuationbox}
        % Take advantage of the already applied Pygments mark-up to insert
        % potential linebreaks for TeX processing.
        %        {, <, #, %, $, ' and ": go to next line.
        %        _, }, ^, &, >, - and ~: stay at end of broken line.
        % Use of \textquotesingle for straight quote.
        \newcommand*\Wrappedbreaksatspecials {%
            \def\PYGZus{\discretionary{\char`\_}{\Wrappedafterbreak}{\char`\_}}%
            \def\PYGZob{\discretionary{}{\Wrappedafterbreak\char`\{}{\char`\{}}%
            \def\PYGZcb{\discretionary{\char`\}}{\Wrappedafterbreak}{\char`\}}}%
            \def\PYGZca{\discretionary{\char`\^}{\Wrappedafterbreak}{\char`\^}}%
            \def\PYGZam{\discretionary{\char`\&}{\Wrappedafterbreak}{\char`\&}}%
            \def\PYGZlt{\discretionary{}{\Wrappedafterbreak\char`\<}{\char`\<}}%
            \def\PYGZgt{\discretionary{\char`\>}{\Wrappedafterbreak}{\char`\>}}%
            \def\PYGZsh{\discretionary{}{\Wrappedafterbreak\char`\#}{\char`\#}}%
            \def\PYGZpc{\discretionary{}{\Wrappedafterbreak\char`\%}{\char`\%}}%
            \def\PYGZdl{\discretionary{}{\Wrappedafterbreak\char`\$}{\char`\$}}%
            \def\PYGZhy{\discretionary{\char`\-}{\Wrappedafterbreak}{\char`\-}}%
            \def\PYGZsq{\discretionary{}{\Wrappedafterbreak\textquotesingle}{\textquotesingle}}%
            \def\PYGZdq{\discretionary{}{\Wrappedafterbreak\char`\"}{\char`\"}}%
            \def\PYGZti{\discretionary{\char`\~}{\Wrappedafterbreak}{\char`\~}}%
        }
        % Some characters . , ; ? ! / are not pygmentized.
        % This macro makes them "active" and they will insert potential linebreaks
        \newcommand*\Wrappedbreaksatpunct {%
            \lccode`\~`\.\lowercase{\def~}{\discretionary{\hbox{\char`\.}}{\Wrappedafterbreak}{\hbox{\char`\.}}}%
            \lccode`\~`\,\lowercase{\def~}{\discretionary{\hbox{\char`\,}}{\Wrappedafterbreak}{\hbox{\char`\,}}}%
            \lccode`\~`\;\lowercase{\def~}{\discretionary{\hbox{\char`\;}}{\Wrappedafterbreak}{\hbox{\char`\;}}}%
            \lccode`\~`\:\lowercase{\def~}{\discretionary{\hbox{\char`\:}}{\Wrappedafterbreak}{\hbox{\char`\:}}}%
            \lccode`\~`\?\lowercase{\def~}{\discretionary{\hbox{\char`\?}}{\Wrappedafterbreak}{\hbox{\char`\?}}}%
            \lccode`\~`\!\lowercase{\def~}{\discretionary{\hbox{\char`\!}}{\Wrappedafterbreak}{\hbox{\char`\!}}}%
            \lccode`\~`\/\lowercase{\def~}{\discretionary{\hbox{\char`\/}}{\Wrappedafterbreak}{\hbox{\char`\/}}}%
            \catcode`\.\active
            \catcode`\,\active
            \catcode`\;\active
            \catcode`\:\active
            \catcode`\?\active
            \catcode`\!\active
            \catcode`\/\active
            \lccode`\~`\~
        }
    \makeatother

    \let\OriginalVerbatim=\Verbatim
    \makeatletter
    \renewcommand{\Verbatim}[1][1]{%
        %\parskip\z@skip
        \sbox\Wrappedcontinuationbox {\Wrappedcontinuationsymbol}%
        \sbox\Wrappedvisiblespacebox {\FV@SetupFont\Wrappedvisiblespace}%
        \def\FancyVerbFormatLine ##1{\hsize\linewidth
            \vtop{\raggedright\hyphenpenalty\z@\exhyphenpenalty\z@
                \doublehyphendemerits\z@\finalhyphendemerits\z@
                \strut ##1\strut}%
        }%
        % If the linebreak is at a space, the latter will be displayed as visible
        % space at end of first line, and a continuation symbol starts next line.
        % Stretch/shrink are however usually zero for typewriter font.
        \def\FV@Space {%
            \nobreak\hskip\z@ plus\fontdimen3\font minus\fontdimen4\font
            \discretionary{\copy\Wrappedvisiblespacebox}{\Wrappedafterbreak}
            {\kern\fontdimen2\font}%
        }%

        % Allow breaks at special characters using \PYG... macros.
        \Wrappedbreaksatspecials
        % Breaks at punctuation characters . , ; ? ! and / need catcode=\active
        \OriginalVerbatim[#1,codes*=\Wrappedbreaksatpunct]%
    }
    \makeatother

    % Exact colors from NB
    \definecolor{incolor}{HTML}{303F9F}
    \definecolor{outcolor}{HTML}{D84315}
    \definecolor{cellborder}{HTML}{CFCFCF}
    \definecolor{cellbackground}{HTML}{F7F7F7}

    % prompt
    \makeatletter
    \newcommand{\boxspacing}{\kern\kvtcb@left@rule\kern\kvtcb@boxsep}
    \makeatother
    \newcommand{\prompt}[4]{
        {\ttfamily\llap{{\color{#2}[#3]:\hspace{3pt}#4}}\vspace{-\baselineskip}}
    }
    

    
    % Prevent overflowing lines due to hard-to-break entities
    \sloppy
    % Setup hyperref package
    \hypersetup{
      breaklinks=true,  % so long urls are correctly broken across lines
      colorlinks=true,
      urlcolor=urlcolor,
      linkcolor=linkcolor,
      citecolor=citecolor,
      }
    % Slightly bigger margins than the latex defaults
    
    \geometry{verbose,tmargin=1in,bmargin=1in,lmargin=1in,rmargin=1in}
    
    \setcounter{secnumdepth}{0}
    
\begin{document}
    
    \maketitle
    
    

    
    \subsection{Problem Set 14}\label{problem-set-14}

\subsubsection{Problem 5.19}\label{problem-5.19}

\subsubsection{Problem 5.20}\label{problem-5.20}

\subsubsection{Problem 5.28}\label{problem-5.28}

\subsubsection{Problem 5.30}\label{problem-5.30}

\begin{center}\rule{0.5\linewidth}{0.5pt}\end{center}

    \subsection{Problem Set 13}\label{problem-set-13}

\subsubsection{Problem 5.1}\label{problem-5.1}

\subsubsection{Problem 5.5}\label{problem-5.5}

\subsubsection{Problem 5.8}\label{problem-5.8}

\subsubsection{Problem 5.11}\label{problem-5.11}

\begin{center}\rule{0.5\linewidth}{0.5pt}\end{center}

    \subsection{Problem Set 12}\label{problem-set-12}

\subsubsection{Problem 4.1}\label{problem-4.1}

\subsubsection{Problem 4.2}\label{problem-4.2}

\subsubsection{Problem 4.10}\label{problem-4.10}

\subsubsection{Problem 4.15}\label{problem-4.15}

\begin{center}\rule{0.5\linewidth}{0.5pt}\end{center}

    \subsection{Problem Set 11}\label{problem-set-11}

\subsubsection{Problem 3.31}\label{problem-3.31}

\subsubsection{Problem 3.32}\label{problem-3.32}

We have already done this one in class.

\begin{center}\rule{0.5\linewidth}{0.5pt}\end{center}

    \subsection{Problem Set 10}\label{problem-set-10}

\subsubsection{Problem 3.24}\label{problem-3.24}

\subsubsection{Problem 3.25 (but not f)}\label{problem-3.25-but-not-f}

\begin{center}\rule{0.5\linewidth}{0.5pt}\end{center}

    \subsection{Problem Set 9}\label{problem-set-9}

\subsubsection{Problem 3.1}\label{problem-3.1}

\subsubsection{Problem 3.10}\label{problem-3.10}

\subsubsection{Problem 3.14}\label{problem-3.14}

\begin{center}\rule{0.5\linewidth}{0.5pt}\end{center}

    \subsection{Problem Set 8}\label{problem-set-8}

\subsubsection{Problem 2.26}\label{problem-2.26}

\subsubsection{Problem 2.32}\label{problem-2.32}

\begin{center}\rule{0.5\linewidth}{0.5pt}\end{center}

    \subsection{Problem Set 7}\label{problem-set-7}

\subsubsection{Problem 2.16}\label{problem-2.16}

\subsubsection{Problem 2.18}\label{problem-2.18}

The ``hint'' in this problem probably needs another hint. Think about
the fact that you can express \(N!\) as \(N\times (N-1)!\), and that
means you can rearrange that formula to write
\[ (N-1)! = \frac{N!}{N} \]

Use a similar technique on \((q+N-1!)\) and then simplify.

Also another hint for later on, don't forget that \(N = N^1 = N^{2/2}\).

\subsubsection{Problem 2.22}\label{problem-2.22}

Your hint for this one is in part (d). The question asks, ``Out of all
the macrostates, what fraction have reasonably large probabilities?''
What he is getting at is you have found in part (c) how tall the
multiplicity peak is, and in part (b) you have found the area under the
multiplicity curve so using the loose approximation of the area as a
thin rectangle the width of that rectangle would correspond to the
number of macrostates that have a reasonable change of occuring. The
``fraction'' that have a reasonable chance of occuring is just this
number divided by the total number of macrostates. This will be a very
small number. This small number represents the fraction of all the
macrostates that are reasonably likely to occur.

\subsubsection{Problem 2.28}\label{problem-2.28}

\subsubsection{Problem 2.29}\label{problem-2.29}

I have already done this in jupyter, so just make sure you have added it
to your personal version of that file and know how it works.

There is a question at the end of this problem that says, ``Also compute
the entropy over long time scales, assuming that all microstates are
accessible.'' What this means is that instead of plugging in the
multiplicity of the of the most likely macrostate into the entropy
formula, you are pluggin in the total multiplicity of all the
macrostates. This is just the sum of the multiplicities of the
macrostates.

Another way you can get this number is to think about the system as a
\emph{single} solid (of 500 particles) with a 100 packets of enery to
share. What is the multiplicity of this single solid. This number is
equal to the sum of the total multiplicities of the two solids in
contact. Show yourself that this is true.

In jupyter to get this to work on the pandas dataframe that we made a
couple of things may be handy. First, make sure your dataframe expresses
all of its values as \emph{float numbers} (which just means scientific
notation) rather than \emph{integers} which is so long your computer can
not display it. To do this run something like this on your table (mine
was called \texttt{table2}).

\begin{verbatim}
table2 = table2.astype(float)
\end{verbatim}

Next, you can sum up the \texttt{multi\_total} column using this
function:

\begin{verbatim}
table2['multi_total'].sum()
\end{verbatim}

Hope this helps.

\begin{center}\rule{0.5\linewidth}{0.5pt}\end{center}

    \subsection{Problem Set 6}\label{problem-set-6}

\subsubsection{Problem 2.1}\label{problem-2.1}

\subsubsection{Problem 2.5 a and b}\label{problem-2.5-a-and-b}

\subsubsection{Problem 2.9}\label{problem-2.9}

We have done the first part of this in Google Sheets and in Jupyter. You
can do the second part in either. Just copy the results of the table and
graph into your homework.

\subsubsection{Problem 2.10}\label{problem-2.10}

This is also easy to do in Jupyter now that we have the code written.

\subsubsection{Problem 2.12 d}\label{problem-2.12-d}

Do this by graph this function \(y=ln(1+x)\) in the vicinity near x = 0
and then use the result from part (c) which is that the derivative of
natural log is the inverse function:

\[ \frac{d \ln(x)}{dx} = \frac{1}{x} \]

Think about the slope of the function and what the equation of a line
that approximates this would be.

\subsubsection{Problem 2.13 b}\label{problem-2.13-b}

I already used this trick in the notes, but I want you identify it and
to write it down.

\begin{center}\rule{0.5\linewidth}{0.5pt}\end{center}

    \subsection{Problem Set 5}\label{problem-set-5}

\subsubsection{Finish problem 1.34}\label{finish-problem-1.34}

You have done the work part of (a), so use the equipartition theorem to
finish the rest. Don't forget it is a diatomic gas.

\subsubsection{Problem 1.36}\label{problem-1.36}

It might help to do 1.35 first but I'll tell you the gist and then you
can just use the results. We derived an expression in class for initial
and final volume and temperature ratios for an adiabatic compression:

\[ \left(\frac{T_f}{T_i}\right)^{f/2}=\frac{V_i}{V_f}\]

This is also the same information as equation 1.38 and 1.39 from the
book. Use your proportionalities from the ideal gas law to get
\(T\propto PV\), and then plug that in and simplify. You will get
equation 1.40 from the book, as long as you let \(\gamma = (f+2)/f\).
Use this form to help with Problem 1.36.

\subsubsection{Problem 1.38}\label{problem-1.38}

Great conceptual question. Think about which bubble with expand more,
the one that experienced isothermal expansion, or the one that
experienced adiabatic expansion. They experience the same change in
pressure since they start at the same level in the water and both reach
the surface.

\begin{center}\rule{0.5\linewidth}{0.5pt}\end{center}

    \subsection{Problem Set 4}\label{problem-set-4}

\subsubsection{Problem 1.7}\label{problem-1.7}

\subsubsection{Problem 1.16}\label{problem-1.16}

\subsubsection{Problem 1.17}\label{problem-1.17}

Additional Hint: you can use pressure in atm and volume in cm\(^3\) if
you use the gas constant \$ R=83.1\$ which has units of
\(\frac{atm \cdot cm^3}{mol K}\).

\begin{center}\rule{0.5\linewidth}{0.5pt}\end{center}

    \subsection{Problem Set 3}\label{problem-set-3}

\subsubsection{Equilibrium Temperature}\label{equilibrium-temperature}

A 5kg block of aluminum (\(c_{Al} = 0.9 kJ/kgK\)) at 500C is in thermal
contact with a 5kg block of lead (\(c_{Pb} = 0.13kJ/kgK\)) at 75C. What
is their equilibrium temperature?

\subsubsection{Problem 1.41}\label{problem-1.41}

\subsubsection{Problem 1.47 or 1.48 your
choice}\label{problem-1.47-or-1.48-your-choice}

\begin{center}\rule{0.5\linewidth}{0.5pt}\end{center}

    \subsection{Problem Set 2}\label{problem-set-2}

\subsubsection{Problem 1.28}\label{problem-1.28}

Estimate how long it should take to bring a cup of water to boiling
temperature in a typical 600-watt microwave oven, assuming that all the
energy ends up in the water. (Assume any reasonable initial temperature
for the water.) Explain why no heat is involved in this process.

\subsubsection{Problem 1.33 a}\label{problem-1.33-a}

We did this in class already, but just make sure you have it written
down, so we can come back and finish the rest later.

\subsubsection{Problem 1.34 a}\label{problem-1.34-a}

But ONLY the work done on each stage of the cycle. We'll come back and
talk about the heat added and the change in th eenergy content later.
Express your answers in terms of \(P_1, P_2, V_1, V_2.\)

\begin{center}\rule{0.5\linewidth}{0.5pt}\end{center}

    \subsection{Problem Set 1}\label{problem-set-1}

\subsubsection{Problem 1.1}\label{problem-1.1}

\subsubsection{Problem 1.3}\label{problem-1.3}

\subsubsection{Problem 1.12}\label{problem-1.12}

\subsubsection{Proportionality}\label{proportionality}

Proportionality is a very convenient way to summarize a relationship,
but its the kind of thing that everyone assumes someone else taught you.

So when we see an equation like the ideal gas law \(P V = N k_B T\) this
equation contains many relationships that can be stated as
\emph{ratios}. As an example, imagine we are conducting an experiment
where we vary the temperature of a gas, but we do not change the number
of particles or the volume of the container. The question is what
happens to the pressure? We can relate the \emph{ratio} of the two
temperatures to the \emph{ratio} of the two different pressures. Here is
how:

Solve for the Pressure:

\[ P = \frac{N k_B}{V} T \]

and then see that this general equation must be satisfied in the
specific instances of \(T_1, T_2, P_1, P_2\):
\[ P_1 = \frac{N k_B}{V} T_1 \] \[ P_2 = \frac{N k_B}{V} T_2 \]

Dividing the two equations as 2/1, shows that all of the constant things
(\(N, k_b, V\)) cancel out and we are left with a proportion

\[\frac{P_2}{P_1} = \frac{T_2}{T_1}\]

This type of proportionality is called \emph{directly proportional} or
\emph{linearly proportional} since if you double the temperature
\(T_2/T_1 = 2\) then the equation says that the pressure would double as
well \(P_2/P_1 = 2\).

The fancy/mathy way of writing this is \(P \propto T\). When you see
that statement, it means you can write down
\[\frac{P_2}{P_1} = \frac{T_2}{T_1}\]

Other types of proportionality follow from this logic. So for example

\[ y \propto \frac{1}{x} \]

is called \emph{inverse proportionality}. It is also sometimes written
as \(y \propto x^{-1}\). From these statements, you can write down

\[ \frac{y_2}{y_1} = \frac{x_1}{x_2} = \left(\frac{x_2}{x_1}\right)^{-1} \]

There are other proportionality statements that are very important in
physics like the famous \emph{inverse square law} which comes in the
form \(y \propto x^{-2}\). These kinds of proportionality statements can
also be written as an equation with a \emph{constant of proportionality}
which is often the mores familiar form to us. So for example the law of
Universal Gravitation can be written like

\[ F = \frac{G m_1 m_2}{r^2} \]

We could say the the force is directly proportional to each mass, and
inversely proportional to the square of the distance between them. The
constant of proportionality is \emph{G}.

Ok, now go back to the ideal gas law and imagine the following
scenarios:

\begin{enumerate}
\def\labelenumi{\arabic{enumi}.}
\tightlist
\item
  By what factor does the pressure change if the temperature triples
  (here we assume that all other things are unchanged)? (that language
  of ``by what factor'' is just another way of saying ``whats the ratio
  of pressures'')
\item
  By what factors does the pressure change if the number of particles
  doubles? What about halving? (remember that this is only true if the
  temperature is the same as well as the volume - how would you do
  that?).
\item
  By what factor does the pressure change if the volume doubles?
\item
  By what factor does the pressure change if the volume doubles and the
  temperature triples?
\item
  By what factors would the temperature change if the volume doubled and
  the pressure tripled?
\item
  By what factors would the volume change if the number of particles
  halved and the pressure doubled and the temperature \emph{decreased}
  by a factor of 10? How would you do this in the lab?
\end{enumerate}


    % Add a bibliography block to the postdoc
    
    
    
\end{document}
